% Part: incompleteness
% Chapter: theories-computability
% Section: inseparability

\documentclass[../../../include/open-logic-section]{subfiles}

\begin{document}

\olfileid{inc}{tcp}{ins}

\olsection{په~$\Th{Q}$ کښې ثابتېدونکې او ردېدونکې جملې په محاسبوي
  ډول نه بېلېدونکې دي}


پرېږدئ~$\Th{\bar Q}$ د هغو جملو سټ وي چې \emph{نفيانې} ئې په
$\Th{Q}$ کښې ثابتېدونکې دي؛ يعنې
$\Th{\bar Q} = \Setabs{!A}{\Th{Q} \Proves \lnot !A}$۔ په ياد ولرئ
چې بېل سټونه~$A$ او~$B$ هله په محاسبوي ډول نه بېلېدونکي بلل کېږي
چې داسې محاسبه کېدونکے سټ~$C$ نۀ وي چې~$A \subseteq C$ او
$B \subseteq \Complement{C}$۔

\begin{lem}
$\Th{Q}$ او~$\Th{\bar Q}$ په محاسبوي ډول نه بېلېدونکي دي۔
\end{lem}

\begin{proof}
فرض کړئ~$C$ داسې محاسبه کېدونکے سټ دے چې~$\Th{Q} \subseteq C$ او
$\Th{\bar Q} \subseteq \Complement{C}$۔ پرېږدئ~$R(x,y)$ لاندې
اړيکه وي:
\begin{quote}
$x$ د~$!D(u)$ فارمول کوډوي او~$!D(\num y)$ په~$C$ کښې دے۔
\end{quote}
موږ به وښيو چې~$R(x,y)$ يوه نړيواله محاسبه کېدونکې اړيکه ده، چې
تناقض ترې ترلاسه کېږي۔

فرض کړئ~$S(y)$ محاسبه کېدونکې ده، او په~$\Th{Q}$ کښې د~$!D_S(u)$
په وسيله تمثيلېږي۔ بيا
\begin{eqnarray*}
S(n) & \lif & \Th{Q} \Proves !D_S(\num n) \\
& \lif & !D_S(\num n) \in C
\end{eqnarray*}
او
\begin{eqnarray*}
\lnot S(n) & \lif & \Th{Q} \Proves \lnot !D_S(\num n) \\
& \lif & !D_S(\num n) \in \Th{\bar Q} \\
& \lif & !D_S(\num n) \not\in C
\end{eqnarray*}
\textbf{د ژباړې د نسخې سمون (OLCMP-094):} سرچينې په دواړو
ماوراءژبنیو اړيکو کښې طبيعي عدد د څيز-ژبې د عددنښې په بڼه ليکلے و؛
دلته عددنښه يوازې د تمثيلوونکي فارمول دننه کارول شوې ده۔
نو~$S(y)$ له~$R(\Gn{!D_S(u)},y)$ سره هم‌ارزه ده۔
\textbf{د ژباړې د نسخې سمون (OLCMP-095):} سرچينې په وروستۍ کرښه
کښې د فارمول کوډ په داسې بڼه ليکلے و چې ازاد متغير ئې په نامعتبره
عددنښه بدلوله او د څپرکي د ګوډل-کوډ نښه ئې نۀ کاروله؛ دلته د هماغه
يو-ازاد-متغير فارمول سم کوډ راغلے دے۔
\end{proof}

\end{document}
